\documentclass[10pt,a4paper]{amsart}
\usepackage[margin=19mm]{geometry}
\usepackage{amsmath,amssymb,mathtools}
\usepackage[scr=rsfs,cal=euler]{mathalfa}
\usepackage{xfrac}
\usepackage[unicode,hidelinks,bookmarks=false]{hyperref}
\newcommand{\ie}{i.e.\ }
\newcommand{\cf}{cf.\ }
\newcommand{\resp}{respectively\ }
\newcommand{\Subsection}{\subsection}
\providecommand{\nobreakdash}{\nobreak\hbox{-}}
\setlength{\emergencystretch}{2em}
\multlinegap0pt
\usepackage{amssymb}
\usepackage[normalem]{ulem}
\usepackage[shortlabels]{enumitem}
\newtheorem{definition}{Definition}
\newtheorem{proposition}{Proposition}
\usepackage{cleveref}
\crefname{definition}{Definition}{Definitions}
\crefname{proposition}{Proposition}{Propositions}
\newcommand{\N}{\mathbb N}
\newcommand{\makeset}[2]{\left\lbrace #1 \;\middle|\;
 \begin{tabular}{@{}l@{}}
   #2
  \end{tabular}
  \right\rbrace}
\title{Maximal subgroups of homeomorphism groups}
\author{S. Bardyla, L. Elliott,, Y. P\'eresse}
\date{Source: arXiv:2608.28211v1 (CC BY 4.0)}
\begin{document}
\maketitle

\noindent\emph{Source and licence.} Maximal subgroups of homeomorphism groups. Version arXiv:2608.28211v1 (CC BY 4.0), distributed by arXiv under the Creative Commons Attribution 4.0 International licence, http://creativecommons.org/licenses/by/4.0/, as recorded in the arXiv licence metadata for this exact version and shown on the abstract page https://arxiv.org/abs/2608.28211v1. Full licence notice: \texttt{licenses/arxiv-2608.28211-CC-BY-4.0.txt}.\par\smallskip

\noindent\textit{Editorial scope.} Counted unit: the article's proposition kappa (source0.tex lines 401--404). The article's complete original proof of it is delivered with it (source0.tex lines 405--436, one \texttt{proof} environment of 32 lines). No mathematics is written, repaired or re-derived: the statement and the proof are the article's own lines, in the article's own notation, and no retained line is substituted or re-flowed.  Retained uncounted as the source-local context the proof needs: lines 199--207; lines 211--213.  Nothing was omitted from any retained span, so the package carries no omission record.  No retained line contains a citation, so the wrapper carries no bibliography and no bibliography entry.  No retained line contains a figure environment, an image-inclusion command or drawing code, so no image is delivered and the package declares no asset dependency.  Editorial formatting only, in addition: an AMS \texttt{amsart} preamble that restates the article's own declarations and macros used by the retained lines, namely the environments \texttt{definition}, \texttt{proposition} on separate counters, together with the standard packages those lines need.  The retained environments print in this document as Definition 1, Proposition 1 in order of appearance, whereas the article prints the same items with the numbering of its own full document.  The complete native source of the article as distributed in the arXiv e-print for this exact version is bundled unchanged as \texttt{sources/208768/source0.tex}.
\begin{definition}\label{clastic}
    A non-empty topological space \(X\) is called \emph{homeoclastic} if \(X\) satisfies the following conditions:
    \begin{enumerate}
        \item \(X\) is Hausdorff;
        \item the topological moieties of \(X\) form a subbasis for \(X\);
         \item whenever \(\{A,B\}\) is a partition of \(X\) into two clopen sets, then either \(A\) or \(B\) is homeomorphic to \(X\);
        \item a clopen subspace of \(X\) is homeomorphic to \(X\) if and only if it contains a topological moiety of \(X\).
    \end{enumerate}
\end{definition}

\begin{definition}\label{def:kappaclast}
  For an infinite cardinal $\kappa$ a space $X$ is called {\em $\kappa$-homeoclastic} if $X$ is homeoclastic and $X$  can be partitioned into \(\kappa\) topological moieties.
\end{definition}

\begin{proposition}\label{kappa}
    Suppose that \(\kappa\) is an infinite cardinal with the discrete topology and \(X\) is an infinite Hausdorff topological space with a basis of clopen sets such that every nonempty clopen subspace of \(X\) is homeomorphic to \(X\).
    If for all cardinals \(\gamma <\kappa\) with the discrete topology we have that \(\gamma \times X \) is not homeomorphic to \(\kappa \times X\), then \(\kappa \times X\) is a \(\kappa\)-homeoclastic space.
\end{proposition}

\begin{proof}
    First we check the four conditions from Definition~\ref{clastic}.
    
    (1) The product space $\kappa \times X$ is Hausdorff since the discrete space $\kappa$ and $X$ are both Hausdorff.

    (2) Let \(B\) be the set of topological moieties of \(X\). Since each non-empty clopen subset of $X$ is homeomorphic to $X$, it follows that $B$ coincides with the family of all proper non-empty clopen subsets of $X$. Since $X$ has a basis of clopen sets, the sets \(\{\alpha\}\times U\) where \(\alpha\in \kappa\) and \(U\in B\) form a basis for \(\kappa \times X\).
    It follows that the family $$S=\makeset{\bigcup_{\alpha\in M}(\{\alpha\}\times U_\alpha)}{\(M \hbox{ is a moiety of }\kappa \hbox{ and } U_\alpha\in B \hbox{ for all }\alpha\in M\)}$$ forms a subbase of $\kappa \times X$. Notice that $S$ consists of topological moieties of $\kappa\times X$. 

    (3) Suppose that \(\{A,B\}\) is a partition of \(\kappa \times X \) into clopen sets. Without loss of generality, suppose that \(I:=\makeset{\alpha\in \kappa}{\((\{\alpha\}\times X)\cap A \neq \emptyset\)}\) has cardinality \(\kappa\). For each \(i\in I\), we have that \((\{i\}\times X)\cap A\) is homeomorphic to a non-empty clopen subspace of \(X\). By the assumption, each of the spaces \((\{i\}\times X)\cap A\)  for \(i\in I\) is homeomorphic to \(X\). As \(A\) is the topological sum of the spaces \((\{i\}\times X)\cap A\) for \(i\in I\) and $|I|=\kappa$, it follows that \(A\) is homeomorphic to \(\kappa \times X\).
    
   %  So \(A\) is homeomorphic to the topological sum of \(\kappa \times X\) with a clopen subset $U$ of \(\kappa \times X\). For each $\alpha\in\N$ let $U_\alpha=\{x\in X| (\alpha,x)\in U\}$. Thus \(A\) is homeomorphic to the topological sum over \(\alpha\in \N\) of the topological sum of the spaces \(X\times\{0\}\) and \(U_\alpha{\times}\{1\}\).
   % As each \(U_\alpha\) is a clopen subspace of \(X\), we either have \(U_\alpha\) homeomorphic to \(X\) or \(X\backslash U_\alpha\) homeomorphic to \(X\). 
   % Since $X$ is homeoclastic, there exists a partition $\{C,D\}$ of $X$ such that $C\cong D\cong X$. If $U_\alpha\cong X$, then $$X=C\cup D\cong (X\times \{0\})\cup (U_\alpha\times\{1\}).$$ If $X\setminus U_\alpha\cong X$, then $$X=(X\setminus U_\alpha)\cup U_\alpha\cong (X\times\{0\})\cup  (U_\alpha\times\{1\}),$$ as $U_\alpha$ is clopen in $X$. It follows that \(A\) is homeomorphic to \(\kappa \times X\).

    %Suppose that \(U\) is a clopen subset of \(\kappa \times X\). If \(U\) contains a TM, then \(U\) is homeomorphic to the topological sum of \(\kappa\times X\) with a clopen subspace of \(\kappa\times X\). So in particular there are clopen subspaces \((U_{\alpha})_{\alpha \in \kappa}\) of \(X\) such that
    %\[U\cong (\sqcup_{\alpha\in \kappa} X)\sqcup (\sqcup_{\alpha\in \kappa} U_{\alpha}) \cong \sqcup_{\alpha\in \kappa} (X\sqcup U_{\alpha})\cong \sqcup_{\alpha\in \kappa} X\cong \kappa \times X.\]
    %We can express this sum as the topological sum of \(\kappa\) spaces each of which are topological sum of \(X\) with a clopen subspace of \(X\). 
    %As \(X\) is homeomorphic to all non-empty clopen subspaces of the topological sum of \(X\) with itself, it follows that \(U\) is homeomorphic to \(\kappa \times X\) as required.

(4) Suppose that \(U\) is a clopen subset of \(\kappa \times X\). If \(U\) contains a topological moiety $T$, then $T$ is homeomorphic to $\kappa \times X$ and $V=U \setminus T$ is clopen in $\kappa \times X$. Hence $V$ is a union $V=\bigcup_{\alpha \in \kappa} \{\alpha\} \times V_{\alpha}$ where each $V_{\alpha}$ is clopen (possibly empty) in $X$. If we use $\sqcup$ to denote topological sums, then 
    \[U=T \sqcup V\cong (\sqcup_{\alpha\in \kappa} X)\sqcup (\sqcup_{\alpha\in \kappa} V_{\alpha}) \cong \sqcup_{\alpha\in \kappa} (X\sqcup V_{\alpha})\cong \sqcup_{\alpha\in \kappa} X\cong \kappa \times X.\]
    
    Suppose now that \(U\) is homeomorphic to \(\kappa\times X\). We need to show that \(U\) contains a topological moiety of \(\kappa\times X\).
   For each $\xi\in\kappa$ the set \(U_{\xi}:=U\cap (\{\xi\}\times X)\) is either empty or homeomorphic to \(X\) since it is clopen in \(\{\xi\}\times X \cong X\). 
   Let \(I_U:=\makeset{\xi\in \kappa}{\(U_\xi\neq \emptyset\)}\). 
   Note that \(U\) is homeomorphic to \(I_U \times X\).
   So if \(|I_U|< \kappa\) then, by the assumption, \(U\) is not homeomorphic to \(\kappa \times X\), a contradiction. Thus \(|I_U|=\kappa\). Let \(M\) be a moiety of \(I_U\). 
   Then \(\bigcup_{\xi\in M} U_\xi\) is a topological moiety of \(\kappa \times X\) contained in \(U\), as required.

   
    To verify that the condition of \cref{def:kappaclast} holds, let \(\makeset{M_\alpha}{$\alpha< \kappa$}\) be a partition of \(\kappa\) into moieties. Then \(\makeset{M_\alpha \times X}{$\alpha< \kappa$}\) is a partition of \(\kappa\times X\) into topological moieties. Hence \(\kappa \times X\) is a \(\kappa\)-homeoclastic space.
\end{proof}

\end{document}
